% mathdown core: what the LaTeX produced by filters/mathdown.lua needs.
% Included by defaults/mathdown.yaml. Change styling in a style file
% (such as mathdown-default.tex), not here.

\usepackage{mathtools}   % \coloneqq for :=
\usepackage{amsthm}      % proof, \qedhere, \newtheorem

% Unicode symbols in math (⊆, ⟹, ±, …): unicode-math-input for pdflatex;
% with xelatex/lualatex pandoc's template already loads unicode-math.
\usepackage{iftex}
\ifPDFTeX\usepackage{unicode-math-input}\fi

% The environments the filter writes for **Theorem**, **Definition**, …
% They are defined at \begin{document}, and only if a style or your own
% preamble has not defined them already, so those can choose their look,
% numbering or names.
\makeatletter
\newcommand\mathdown@thm[3][mdtheorem]{%
  \@ifundefined{#2}{\newtheorem{#2}[#1]{#3}}{}}
\AtBeginDocument{%
  \@ifundefined{mdtheorem}{\theoremstyle{plain}\newtheorem{mdtheorem}{Theorem}}{}%
  \theoremstyle{plain}%
  \mathdown@thm{mdlemma}{Lemma}%
  \mathdown@thm{mdproposition}{Proposition}%
  \mathdown@thm{mdcorollary}{Corollary}%
  \mathdown@thm{mdconjecture}{Conjecture}%
  \mathdown@thm{mdclaim}{Claim}%
  \theoremstyle{definition}%
  \mathdown@thm{mddefinition}{Definition}%
  \mathdown@thm{mdexample}{Example}%
  \mathdown@thm{mdaxiom}{Axiom}%
  \mathdown@thm{mdnotation}{Notation}%
  \theoremstyle{remark}%
  \mathdown@thm{mdremark}{Remark}%
}
\makeatother
